% =====================================================================
% PREAMBLE.TEX
% University of Bologna — Thesis preamble
% All packages, layout and formatting configuration live here.
% main.tex only orchestrates \input{} calls and stays clean.
% =====================================================================

% ---------------------------------------------------------------------
% Encoding & Language
% ---------------------------------------------------------------------
\usepackage[T1]{fontenc}
\usepackage[utf8]{inputenc}
\usepackage[english]{babel}

% ---------------------------------------------------------------------
% Page Layout
% ---------------------------------------------------------------------
\usepackage{geometry}
\geometry{
	left=25mm,
	right=25mm,
	top=25mm,
	bottom=25mm
}

\usepackage{setspace}
\onehalfspacing % use \doublespacing if your program requires it

\usepackage{indentfirst}
\usepackage{pdflscape}

% ---------------------------------------------------------------------
% Fonts
% ---------------------------------------------------------------------
\usepackage{lmodern}        % Improved Computer Modern
\usepackage{microtype}      % Better typography

% ---------------------------------------------------------------------
% Mathematics
% ---------------------------------------------------------------------
\usepackage{amsmath}
\usepackage{amssymb}
\usepackage{amsfonts}
\usepackage{amsthm}
\usepackage{bm}             % Bold math symbols
\usepackage{mathtools}

% ---------------------------------------------------------------------
% Graphics & Figures
% ---------------------------------------------------------------------
\usepackage{graphicx}
\usepackage{float}
\usepackage{subcaption}
\usepackage{caption}
\usepackage{tikz}
\usetikzlibrary{arrows.meta,positioning,calc}

% ---------------------------------------------------------------------
% Tables
% ---------------------------------------------------------------------
\usepackage{booktabs}
\usepackage{multirow}
\usepackage{array}
\usepackage{longtable}
\usepackage{tabularx}

% ---------------------------------------------------------------------
% Algorithms & Code
% ---------------------------------------------------------------------
\usepackage{algorithm}
\usepackage{algpseudocode}

\usepackage{listings}
\usepackage{xcolor}

\lstset{
	basicstyle=\ttfamily\small,
	keywordstyle=\color{blue},
	commentstyle=\color{gray},
	stringstyle=\color{red},
	numbers=left,
	numberstyle=\tiny,
	frame=single,
	breaklines=true
}

% ---------------------------------------------------------------------
% Units & Symbols
% ---------------------------------------------------------------------
\usepackage{siunitx}
\sisetup{
	per-mode=symbol,
	detect-all
}

% ---------------------------------------------------------------------
% Hyperlinks
% (loaded late, but before titlesec/fancyhdr as is standard practice)
% ---------------------------------------------------------------------
\usepackage{hyperref}
\hypersetup{
	colorlinks=true,
	linkcolor=blue,
	citecolor=blue,
	urlcolor=blue
}

% ---------------------------------------------------------------------
% Appendices
% ---------------------------------------------------------------------
\usepackage{appendix}

% ---------------------------------------------------------------------
% Nomenclature / Acronyms
% ---------------------------------------------------------------------
\usepackage{nomencl}
\makenomenclature

% Alternative: glossaries-based acronym list (uncomment if you prefer it
% over nomencl; do not use both systems together)
% \usepackage[acronym]{glossaries}
% \makeglossaries

% ---------------------------------------------------------------------
% PDF Inclusion
% ---------------------------------------------------------------------
\usepackage{pdfpages}

% ---------------------------------------------------------------------
% Line Breaking & Widow Control
% ---------------------------------------------------------------------
\usepackage[all]{nowidow}

% =====================================================================
% Custom Commands & Thesis Metadata
% Edit these three lines for your own thesis — they are reused
% throughout the running headers/footers below.
% =====================================================================
\newcommand{\R}{\mathbb{R}}
\newcommand{\vect}[1]{\boldsymbol{#1}}

\newcommand{\AuthorName}{\textit{X.Y.Surname}}
\newcommand{\ThesisTitle}{Title of Thesis}
\newcommand{\CourseUniversity}{MSc, Statistical Sciences, Department of Statistical Sciences "Paolo Fortunati", University of Bologna}

% =====================================================================
% Chapter & Section Appearance
% =====================================================================
\usepackage{titlesec}

% Chapters: small "Chapter X", big title, better spacing
\titleformat{\chapter}[display]
{\normalfont\bfseries}
{\Large\chaptername\ \thechapter}
{5ex}
{\huge}

% =====================================================================
% Headings & TOC / Running Headers-Footers
% =====================================================================
\usepackage{chngcntr} % for \counterwithin
\usepackage{tocloft}
\usepackage{fancyhdr}
\pagestyle{fancy}
\fancyhf{}
\setlength{\headheight}{14pt}

% -----------------------
% Marks (chapter/section)
% -----------------------
\renewcommand{\chaptermark}[1]{\markboth{\thechapter.\ #1}{}}
\renewcommand{\sectionmark}[1]{\markright{\thesection.\ #1}}

% -----------------------
% Default header (used by the "fancy" style)
% -----------------------
\fancyhead[L]{\small\nouppercase{\leftmark}}  % current chapter/section
\fancyhead[R]{\small\thepage}                 % page number

% -----------------------
% Default footer (used by the "fancy" style)
% -----------------------
% Odd pages -> thesis title on the left
\fancyfoot[L]{%
	\ifodd\value{page}
	\footnotesize \textit{\ThesisTitle}
	\fi
}
% Even pages -> course/university on the right
\fancyfoot[R]{%
	\ifodd\value{page}
	\else
	\footnotesize \CourseUniversity
	\fi
}

\renewcommand{\headrulewidth}{0.4pt}
\renewcommand{\footrulewidth}{0.4pt}

% -----------------------
% "plain": bare style used e.g. for the bibliography's first page
% -----------------------
\fancypagestyle{plain}{
	\fancyhf{}
	\fancyfoot[C]{\thepage}
	\renewcommand{\headrulewidth}{0pt}
	\renewcommand{\footrulewidth}{0pt}
}

% -----------------------
% "chapterstar": first page of unnumbered chapters
% (Table of Contents, List of Figures/Tables, and any future
% Abstract/Acknowledgment/Nomenclature pages)
% -----------------------
\fancypagestyle{chapterstar}{
	\fancyhf{}
	\fancyfoot[C]{\thepage}
	\renewcommand{\headrulewidth}{0pt}
	\renewcommand{\footrulewidth}{0pt}
}

% -----------------------
% "chapterplain": first page of numbered chapters (no header, footer only)
% -----------------------
\fancypagestyle{chapterplain}{
	\fancyhf{}
	\fancyfoot[L]{\footnotesize \CourseUniversity}
	\fancyfoot[R]{\footnotesize \thepage}
	\renewcommand{\headrulewidth}{0pt}
	\renewcommand{\footrulewidth}{0.4pt}
}

% -----------------------
% "chapternum": regular pages within numbered chapters
% (chapter-relative page numbers, e.g. 2.1, 2.2, ...)
% -----------------------
\fancypagestyle{chapternum}{
	\fancyhf{}
	\fancyhead[L]{\small\rightmark}
	\fancyhead[R]{\small\thepage}
	\fancyfoot[R]{%
		\ifodd\value{page}
		\footnotesize \textit{\ThesisTitle}
		\fi
	}
	\fancyfoot[L]{%
		\ifodd\value{page}
		\else
		\footnotesize \CourseUniversity
		\fi
	}
	\renewcommand{\headrulewidth}{0.4pt}
	\renewcommand{\footrulewidth}{0.4pt}
}

% =====================================================================
% End of Preamble
% =====================================================================